\section{Scope, Method, and Qualifications}

\subsection*{Subject and Intended Reader}

Part II of this four-part Catholic series on the Four Last Things (Death, Judgment, Hell, and Heaven) addresses serious Catholic believers seeking a sober, source-grounded exposition of the Church's doctrine on Particular Judgment, General Judgment, and Purgatory, moving from Sacred Scripture and conciliar definitions through the Fathers, Scholastic Doctors, liturgy, and saints into examination of conscience and received prayer. Questions belonging primarily to physical death and preparation for dying (Part I), eternal damnation (Part III), or the beatific vision and the renewed creation (Part IV) are touched only insofar as they intersect directly with Christ's judgment and the final purification of the elect.

\subsection*{Scripture, Liturgical Texts, and Received Prayers}

All English quotations of Sacred Scripture follow the Challoner revision of the Douay-Rheims Bible, paired where helpful with the Clementine Vulgate Latin. Book and verse numbering follows the Vulgate and Douay-Rheims convention throughout---notably 2 Kings 12:13--14 for Nathan's words to David, 2 Machabees 12:43--46 for Judas Machabeus's offering for the dead, Psalm 65:10--12 for the passage through fire and water, Psalm 84:11 for mercy and truth meeting, Psalm 100:1 for mercy and judgment, Psalm 129 for the \textit{De profundis}, Psalm 142:2 for ``enter not into judgment,'' Ecclesiasticus 11:28--29, and Apocalypse 20:11--15.

Latin liturgical texts---the \textit{Memento etiam} and \textit{Nobis quoque peccatoribus} of the Roman Canon, and the Introit, Collect (\textit{Fidelium, Deus}), Gradual, Tract (\textit{Absolve, Domine}), and Sequence (\textit{Dies irae}, strophes 1--8) of the First Mass of the Commemoration of All the Faithful Departed---follow the 1962 \textit{Missale Romanum}, corroborated in the Roman Canon against the 1604 typical edition and in the closing versicles against the 1915 Benziger \textit{Office of the Dead}.

Every text presented for vocal prayer in Section 10 reproduces a received public-domain edition verbatim: the Latin of \textit{Dies irae} strophes 1--8 from the 1962 \textit{Missale Romanum}, its English verse by W.\ J.\ Irons (stanzas 1--4) from the Complete Edition of \textit{Hymns Ancient and Modern} (London: William Clowes and Sons, Limited, catalogued 1890), the Collect \textit{Fidelium, Deus} and closing versicles from the 1962 \textit{Missale Romanum} and the 1915 Benziger \textit{Little Office of the Blessed Virgin Mary and Office of the Dead} (corroborated against the 1889 \textit{Manual of Prayers}), and St.\ Alphonsus Maria de Liguori's prayer on judgment from the 1843 Dublin Duffy edition of \textit{Preparation for Death} (resolving only obvious nineteenth-century OCR spacing glitches without altering wording). Translations of brief Latin liturgical or conciliar phrases embedded inside analytical prose serve solely as explanatory renderings and are not offered for liturgical recitation.

\subsection*{Source Hierarchy, Editions, and Limits}

The exposition distinguishes four principal grades of authority, summarized in Appendix A: Sacred Scripture and solemn dogmatic definitions (\textit{de fide}); ordinary magisterial teaching and Catholic doctrine; patristic and scholastic consensus (\textit{sententia communis}); and open theological questions on which approved Catholic Doctors legitimately differ. Particular care is taken to preserve the doctrinal restraint of the Council of Florence (\textit{Laetentur caeli}, DS 1304) and the Council of Trent (Session XXV, DS 1820), which defined the existence of Purgatory and the efficacy of suffrages while refraining from defining the physical nature of purgatorial fire, the duration of purgation, or the outward mode and location of the General Judgment.

Patristic, scholastic, catechetical, and spiritual texts quoted verbatim come exclusively from identified public-domain editions directly inspected in full: Philip Schaff's \textit{Nicene and Post-Nicene Fathers} (St.\ Augustine's \textit{Confessions}, St.\ John Chrysostom's \textit{Homilies on First Corinthians} and \textit{Philippians}, and St.\ Cyril of Jerusalem's \textit{Catechetical Lectures}), Marcus Dods's 1871 translation of St.\ Augustine's \textit{The City of God} (with E.\ Hoffmann's CSEL 40 Latin), the Oxford Library of the Fathers \textit{Morals on the Book of Job}, the English Dominican translation of St.\ Thomas Aquinas's \textit{Summa Theologiae} and of the posthumously compiled \textit{Supplement} (drawn from St.\ Thomas's \textit{Scriptum super Sententiis}), J.\ O'Sullivan's 1866 translation of St.\ Robert Bellarmine's \textit{Commentary on the Book of Psalms}, J.\ Donovan's 1829 translation of the \textit{Catechism of the Council of Trent}, the 1921 Benziger \textit{Baltimore Catechism No.\ 3}, the Frederick Pustet edition of St.\ Francis de Sales's \textit{Introduction to a Devout Life}, and the 1843 Duffy \textit{Preparation for Death}. Among conciliar witnesses, Lateran IV's \textit{Firmiter credimus} (DS 801) is quoted from its retained Latin text. Other conciliar, papal, canonical, and spiritual texts---the Profession of Faith of Pope Clement IV (1267) sworn at the Second Council of Lyons (1274, DS 856--858), Pope Benedict XII's \textit{Benedictus Deus} (DS 1000--1002), the Council of Florence's \textit{Laetentur caeli} (DS 1304--1306), Council of Trent Sessions VI, XXII, and XXV (DS 1580, 1743, 1753, 1820, 1835), St.\ Gregory the Great's \textit{Dialogues} IV.39, St.\ Catherine of Genoa's \textit{Treatise on Purgatory} (Chapter II), Pope St.\ Paul VI's \textit{Indulgentiarum doctrina} (1967), the 1983 \textit{Code of Canon Law} (cann.\ 992--997), the \textit{Catechism of the Catholic Church} (1997), the Apostolic Penitentiary's \textit{Enchiridion Indulgentiarum} (4th ed., 1999), and Pope Benedict XVI's \textit{Spe salvi} (2007)---are cited by standard Denzinger-Sch\"{o}nmetzer (DS), canon, norm, or section numbers in doctrinal paraphrase only, without verbatim quotation. Where mutable indulgence discipline is discussed in Section 8 and Appendix A, current universal Latin-Church discipline under \textit{Indulgentiarum doctrina} (effective 1967), the 1983 \textit{Code of Canon Law}, and the fourth edition of the \textit{Enchiridion Indulgentiarum} (1999) is dated explicitly and distinguished from the pre-1967 day-and-year reckoning reflected in the 1921 \textit{Baltimore Catechism No.\ 3}.

\subsection*{Rights and Non-Authority Notice}

Offered as an independent aid for reading, reflection, and prayer, this work is not an official liturgical book, a catechism, a magisterial pronouncement, or a substitute for the guidance of a confessor or competent spiritual director. Author-composed exposition carries the license stated in the closing colophon; Scripture, official magisterial and liturgical texts, received prayers and hymns, and historical translations retain their independent public-domain or third-party rights status.
